% ═══════════════════════════════════════════════════════════════════════
\section{Les transformations affines}
% ═══════════════════════════════════════════════════════════════════════

\begin{frame}{Du plan géométrique au plan des coordonnées}
  \small Jusqu'ici, tout a été dit sans calcul (point de vue
  \emph{synthétique}). On choisit maintenant un repère : chaque point
  devient un couple de nombres, chaque transformation devient une
  formule (point de vue \emph{analytique}).
  \[
    d(P,Q) = \norme{P - Q} = \sqrt{(x_P-x_Q)^2 + (y_P-y_Q)^2}.
  \]
  \begin{Rem}{Remarque 1.4 --- Points ou vecteurs ?}
    \small On confond volontairement point $P$ et vecteur
    $\overrightarrow{OP}$ (même colonne de nombres) --- abus commode et
    universel en infographie. Cela explique pourquoi une translation
    n'est \textbf{pas} une application linéaire.
  \end{Rem}
\end{frame}

\begin{frame}{Définition 1.3 --- Application affine du plan}
  \begin{Enonce}{Définition 1.3 --- Application affine du plan}
    Une \defi{application affine} du plan est $w:\R^2\to\R^2$,
    \[
      w\col{x}{y} = \mat{a}{b}{c}{d}\col{x}{y} + \col{e}{f}
      = A\col{x}{y} + \vec{t}, \qquad a,b,c,d,e,f\in\R.
    \]
    $A$ : \defi{partie linéaire} ; $\vec t$ : \defi{partie translation}.
  \end{Enonce}
  \vspace{4pt}
  Une application affine = une application linéaire \textbf{suivie
  d'une translation}. Cette phrase anodine décrit exactement les
  translations, rotations, réflexions, homothéties, cisaillements ---
  et leurs compositions.
\end{frame}

\begin{frame}{Remarque 1.5 --- Pourquoi « affine » et non « linéaire » ?}
  \begin{Rem}{Remarque 1.5}
    Une application linéaire $T$ vérifie $T(\vec 0) = \vec 0$ : elle
    fixe l'origine. Une translation ne fixe \textbf{aucun} point : les
    translations ne sont donc pas linéaires. On ajoute le terme
    $\vec t$ précisément pour les inclure --- d'où « affine ».
  \end{Rem}
  \vspace{6pt}
  \small (Un tour de passe-passe --- les coordonnées homogènes ---
  permettra de récupérer la linéarité, au prix d'une dimension : voir
  section~5.)

  \vspace{10pt}
  \textbf{Vocabulaire fixé pour tout le chapitre :}
  \begin{itemize}
    \item \defi{application affine} : n'importe quelle fonction de
      cette forme, \emph{sans condition sur $A$} ;
    \item \defi{transformation affine} : une application affine
      \textbf{bijective}, i.e. un élément de $\Aff$.
  \end{itemize}
\end{frame}

\begin{frame}{Proposition 1.2 --- Inversibilité}
  \begin{Enonce}{Proposition 1.2 --- Inversibilité}
    \small Une application affine de partie linéaire $A$, translation
    $\vec t$, est bijective si et seulement si $\dete(A) \neq 0$. Alors
    \[
      w^{-1}\col{x}{y} = A^{-1}\col{x}{y} - A^{-1}\vec{t}.
    \]
  \end{Enonce}
  \begin{AuTableau}
    \small Idée : $\dete(A)\neq 0 \Leftrightarrow A$ inversible
    $\Leftrightarrow$ $A$ injective (algèbre linéaire) ; on vérifie
    directement que la formule proposée est une réciproque.
  \end{AuTableau}
  \vspace{6pt}
  \textbf{Conséquence.} Les applications affines inversibles forment un
  sous-groupe $\Aff$ de $\Transf$ : le \defi{groupe affine} du plan.
\end{frame}

\begin{frame}{Proposition 1.3 --- Ce qu'une transformation affine conserve}
  \begin{Enonce}{Proposition 1.3}
    \footnotesize Soit $w$ une transformation affine. Alors :
    \begin{enumerate}
      \item l'image d'une droite est une droite ;
      \item l'image d'une droite par l'origine passe par $(e,f)$ ;
      \item l'image de deux droites parallèles reste parallèle ;
      \item l'image d'un segment est un segment ; le milieu va sur le milieu ;
      \item plus généralement : $w$ préserve les rapports de longueurs
        \emph{sur une même droite} ;
      \item l'image d'un parallélogramme est un parallélogramme.
    \end{enumerate}
  \end{Enonce}
  \begin{AuTableau}
    \footnotesize Idée : écriture paramétrique d'une droite,
    $P + s\vec u$ ; $w(P+s\vec u) = w(P) + s(A\vec u)$, et $A\vec u \neq
    \vec 0$ car $A$ est inversible.
  \end{AuTableau}
\end{frame}

\begin{frame}{Remarque 1.6 --- Ce qu'elle ne conserve \emph{pas}}
  \begin{Rem}{Remarque 1.6}
    Une transformation affine ne préserve \textbf{en général} ni les
    longueurs, ni les angles, ni la perpendicularité, ni les rapports
    de longueurs entre deux directions différentes. Un carré peut
    devenir un parallélogramme quelconque ; un cercle devient une
    ellipse.
    \vspace{4pt}

    En revanche, elle multiplie \textbf{toutes} les aires par le même
    facteur (à venir).
  \end{Rem}
\end{frame}

\begin{frame}{Visualiser : l'image du carré unité}
  \small Le geste le plus utile du chapitre : pour comprendre une
  transformation, regarder ce qu'elle fait au carré unité
  $[0,1]\times[0,1]$. Les \textbf{colonnes de $A$} sont les images de
  $\vec e_1,\vec e_2$ : le carré devient le parallélogramme qu'elles
  engendrent.
  \FigBox[0.5\textheight]{%
    \begin{tikzpicture}[scale=1.0]
      \begin{scope}
        \repere{-0.6}{2.6}{-0.6}{2.6}
        \fill[fig] (0,0) rectangle (1,1);
        \draw[fleche] (0,0) -- (1,0) node[etiq, below right] {$\vec{e}_1$};
        \draw[fleche] (0,0) -- (0,1) node[etiq, above left] {$\vec{e}_2$};
        \node[etiq] at (1.3,-1.15) {le carré unité};
      \end{scope}
      \begin{scope}[xshift=5.2cm]
        \repere{-0.6}{2.6}{-0.6}{2.6}
        \fill[figim] (0,0) -- (1.5,0.4) -- (2.0,1.7) -- (0.5,1.3) -- cycle;
        \draw[fleche, firebrick3] (0,0) -- (1.5,0.4)
          node[etiq, below right] {$A\vec{e}_1 = \col{a}{c}$};
        \draw[fleche, firebrick3] (0,0) -- (0.5,1.3)
          node[etiq, above left] {$A\vec{e}_2 = \col{b}{d}$};
        \node[etiq] at (1.3,-1.15) {son image par $A$};
      \end{scope}
    \end{tikzpicture}}
  \begin{center}\small Figure 1.4\end{center}
\end{frame}

\begin{frame}{Définition 1.4 --- Orientation}
  \begin{Enonce}{Définition 1.4 --- Orientation d'une base, du plan}
    \small $(\vec u,\vec v)$ est \defi{directe} si l'on passe de $\vec
    u$ à $\vec v$, par le plus petit angle, dans le sens
    trigonométrique ; \defi{indirecte} sinon. $w$
    \defi{préserve l'orientation} si elle envoie toute base directe sur
    une base directe ; elle la \defi{renverse} sinon.
  \end{Enonce}
  \FigBox[0.36\textheight]{%
    \begin{tikzpicture}[scale=1.0]
      \begin{scope}
        \draw[fleche] (0,0) -- (1,0) node[etiq, below right] {$\vec u$};
        \draw[fleche] (0,0) -- (55:1) node[etiq, above] {$\vec v$};
        \draw[sectioncolor, thin, -{Stealth[length=1.6mm]}]
          (0.42,0) arc (0:55:0.42);
        \node[etiq] at (0.2,-0.75) {base directe};
      \end{scope}
      \begin{scope}[xshift=3.4cm]
        \draw[fleche] (0,0) -- (1,0) node[etiq, below right] {$\vec u$};
        \draw[fleche] (0,0) -- (-55:1) node[etiq, below right] {$\vec v$};
        \draw[sectioncolor, thin, -{Stealth[length=1.6mm]}]
          (0.42,0) arc (0:-55:0.42);
        \node[etiq] at (0.2,-1.05) {base indirecte};
      \end{scope}
    \end{tikzpicture}}
\end{frame}

\begin{frame}{Proposition 1.4 --- Aire et déterminant}
  \begin{Enonce}{Proposition 1.4 --- Aire et déterminant}
    Soit $w$ affine de partie linéaire $A$, et $D$ une région d'aire
    $S$. Alors $\Aire(w(D)) = \lvert \dete(A)\rvert \cdot S$. Si
    $\dete(A) \neq 0$ : $w$ préserve l'orientation si $\dete(A)>0$, la
    renverse si $\dete(A)<0$.
  \end{Enonce}
  \small (Admis --- relève du théorème du changement de variables.
  Intuition : vrai pour le carré unité par définition du déterminant,
  donc pour tout carré du quadrillage, donc pour toute région
  raisonnable.)
  \begin{Classe}{Le F est-il à l'endroit ?}
    \small Le signe du déterminant : $\dete(A)>0$, la figure se lit
    dans le même sens ; $\dete(A)<0$, elle est retournée comme dans un
    miroir. D'où l'usage d'un motif en « F » dans les figures : un F
    retourné se repère au premier coup d'œil.
  \end{Classe}
\end{frame}

\begin{frame}{Exemple 1.5 --- Image d'un rectangle}
  \small $w\col{x}{y} = \mat{1}{-1}{2}{0}\col{x}{y}+\col{-1}{1}$,
  $R=\intff{0}{2}\times\intff{0}{1}$. Il suffit de calculer l'image des
  4 sommets (prop.~1.3).
  \FigBox[0.42\textheight]{%
    \begin{tikzpicture}[scale=0.62]
      \begin{scope}
        \repere{-2.6}{2.6}{-0.6}{5.6}
        \fill[fig] (0,0) rectangle (2,1);
        \node[etiq] at (1,0.5) {$R$};
      \end{scope}
      \begin{scope}[xshift=7.6cm]
        \repere{-2.6}{2.6}{-0.6}{5.6}
        \fill[figim] (-1,1) -- (1,5) -- (0,5) -- (-2,1) -- cycle;
        \node[etiq, ptim, label={[etiq]below right:$(-1,1)$}] at (-1,1) {};
        \node[ptim, label={[etiq]right:$(1,5)$}] at (1,5) {};
        \node[etiq] at (-0.55,3) {$w(R)$};
      \end{scope}
    \end{tikzpicture}}
  Vérification par les aires : $\dete(A) = 2$, aire de $R$ = 2, donc
  aire de $w(R) = 4$ (confirmé).
\end{frame}

\begin{frame}{Exemple 1.6 --- Image d'un triangle}
  \small Même $w$. Triangle plein $A=(1,1)$, $B=(0,0)$, $C=(-1,1)$
  (isocèle) $\longmapsto$ triangle plein $(-1,3)$, $(-1,1)$, $(-3,-1)$
  (\textbf{pas} isocèle).

  \vspace{6pt}
  \textbf{Une transformation affine ne respecte pas les longueurs.}
\end{frame}
